\section{Introduction}

The geometry of parameterized quantum states controls how changes in circuit parameters alter the represented state. For pure states, this local geometry is captured by the Fubini--Study metric, equivalently the quantum Fisher information matrix up to convention, and motivates quantum natural-gradient optimization \cite{Braunstein1994,Stokes2020}. This viewpoint is attractive because it replaces Euclidean parameter distance by a state-sensitive notion of distance. It also raises a practical question that is usually left implicit: which part of the state-space geometry is actually visible to the learning interface used by the experiment?

A variational learner rarely observes the full state. It interacts with measurement outcomes, retained observables, marginals, or a small classical feature vector. Consequently, two tangent directions that are distinct in the full quantum metric can become indistinguishable after measurement and readout restriction. Conversely, a low-dimensional readout can be strongly aligned with a small set of task-relevant tangent directions even when it retains only a small fraction of the total state-space information. The distinction between full geometry and accessible geometry therefore matters both statistically and algorithmically.

The present work follows a sequence of increasingly operational geometric questions. First, Fubini--Study metric conditioning was used as a trainability control variable: the Sculpture framework meta-learns circuit initializations that improve the condition number of the full metric \cite{AitHaddouBennai2025Sculpting}. Second, under an isotropic tangent model, fixed-basis readout accessibility was shown to obey a rank law: the expected retained tangent mass is set by the readout rank relative to the score-space dimension \cite{AitHaddou2026ReadoutRank}. Third, the isotropy assumption was removed by introducing a spectral orientation framework in which retained information depends jointly on the rank baseline, the tangent-covariance spectrum, and the relative orientation of the readout subspace \cite{AitHaddou2026SpectralGeometry}. The current paper closes the algorithmic loop: instead of using accessibility only as a diagnostic, we use the measurement-accessible tangent metric itself to precondition optimization.

There is related work on measurement-limited feature visibility from a different direction. Gross and Rieser formulate quantum extreme learning machines through Pauli-transfer matrices and characterize which encoded Pauli features are decodable through the row space selected by the reservoir dynamics and measurement operators \cite{GrossRieser2026}. Their problem is representational: a fixed quantum reservoir transforms an input feature library and a classical readout learns from the resulting features. Our problem is differential: the object being projected is the parameter-score tangent of a trainable circuit, and the projected geometry is subsequently used as an optimization metric. The common linear-algebraic skeleton is a projection onto an observable subspace, but the tangent object, normalization, learning dynamics, and operational question are different.

We introduce the \emph{accessible quantum natural gradient} (AQNG). For a commuting retained readout, AQNG builds a covariance-normalized score Jacobian and uses its pullback Gram matrix as the preconditioner. This construction admits an exact projected-score interpretation: the accessible metric is the Gram matrix of the component of the full classical score that lies in the retained feature span. It is therefore not defined by approximating the full QFIM in matrix norm. Instead, the metric is tied to a physically specified measurement interface.

The empirical study is designed to separate three questions that are often conflated: whether a tangent direction is visible to the measurement, whether the resulting accessible metric is numerically well conditioned, and whether the supervised objective still carries a usable gradient signal. We use rank-matched physical, cross-fitted aligned, and random readouts; compare against full QNG, block QNG, Adam, and SGD; test analytic and finite-shot regimes; and include a number-conserving $U(1)$ control in which accessibility and trainability separate sharply. The central conclusion is deliberately narrower than a universal optimizer claim: accessible geometry can be a useful and interpretable optimization resource, but retained tangent information alone does not determine trainability.

The main contributions are:
\begin{enumerate}
    \item an exact accessible-metric construction for retained commuting readouts, together with normalization, damping, and trust controls suitable for optimization;
    \item a rank-matched experimental design that separates physical orientation, cross-fitted tangent alignment, and random-subspace baselines without changing readout rank;
    \item a frozen multi-dataset benchmark and scaling study showing that AQNG can outperform full-QNG and SGD in specific regimes while remaining competitive rather than uniformly superior to Adam;
    \item a diagnostic decomposition of optimization failure into measurement accessibility, metric conditioning, and task-gradient signal, illustrated by the readout-order and $U(1)$ controls.
\end{enumerate}
